\honest{Quantum Explanations}{Eliezer Yudkowsky}{2008}

Many people believe that quantum mechanics is supposed to be confusing. That is a bad frame of mind for a teacher or a student. I am not a physicist, but my essay on Bayes's theorem brings me ``frequent grateful emails.'' Confusion is in our models, not in the world.

So I will not tell you that no one understands quantum mechanics, ``as Richard Feynman once claimed.'' That was true once, ``but we no longer live in that era.'' \nb{The Stanford Encyclopedia of Philosophy reports ``no consensus among physicists or philosophers of physics'' on what quantum theory tells us about the world. A month after this post, in ``Many Worlds, One Best Guess,'' Yudkowsky wrote: ``There is no known reason for the Born probabilities.''} Nor will I repeat what von Neumann ``is reputed to have said'': ``You don't understand quantum mechanics, you just get used to it.'' \nb{As usually reported, von Neumann said it about mathematics: ``in mathematics you don't understand things. You just get used to them.''}

Quantum mechanics is counterintuitive, and the fault is in your intuitions. I will speak of it as perfectly normal and make fun of the intuitions instead.

The standard introduction follows history: particles, then waves, then particles again, then decades of confusion, ``until it was finally sorted out in the second half of the twentieth century.'' I do not say who sorted it out. \nb{The later posts mean many-worlds, proposed by Hugh Everett in 1957, with the decoherence work of H.-D. Zeh and W. Zurek.} An electron is neither a billiard ball nor a wave; it has to be accepted on its own terms.

The order of discovery is not the order of teaching. Physics does not start with biology, so why start with the results of experiments? I grant that stressing experiments ``sounds nice'' to a rationalist. But ``The result of the standard approach is standard confusion.'' I give no evidence about students.

``The classical world is strictly implicit in the quantum world.'' \nb{How single outcomes arise is the measurement problem. The Stanford Encyclopedia says that decoherence does not solve it unless combined with a further account of the theory, such as Everett's, Bohm's or a collapse theory, and there is no consensus on which account is right.} I take a ``strictly realist perspective'': our equations describe the territory. I ask non-realists to wait for a later essay. \nb{The essay came a month later, ``Quantum Non-Realism.'' Pilot-wave theory, in which particles follow definite paths, is realist too; no post in the book discusses it.}

A ``sizable community of scientists'' disputes realism, so you are getting my view. It is ``probably the majority view among theoretical physicists.'' \nb{No evidence is given. In a 2011 poll at a conference on quantum foundations, 24 per cent held that the quantum state is ontic; in a 2025 \textsc{Nature} survey of over 1,100 researchers, about 36 per cent said the wavefunction is real.}

My goal is to make you a native of a quantum universe, not a reluctant tourist. ``Embrace reality. Hug it tight.''
